\QHeading{第11题　解决问题。}
\par\noindent


\Problem{① 亮亮买了一本《格林童话》花了8.55元，又买了一本《小小少年》花了9.5元。买这两本书一共花了多少元？}
\par


\MethodTitle{方法一}
\Answer{{\Formula{\begin{aligned}[t]&\phantom{={}}8.55+9.5\\&=8.55+9.50\\&=18.05\text{（元）}\end{aligned}}}\par 答：一共花了18.05元。}

\MethodTitle{方法二：把9.5元凑成10元}
\Answer{先加10元，多加0.5元再减去。}
\Answer{\Formula{\begin{aligned}[t]&\phantom{={}}8.55+9.5\\
&=8.55+10-0.5\\
&=18.55-0.5\\
&=18.05\text{（元）}\end{aligned}}}
\Answer{答：两本书一共花了18.05元。}
\Problem{② 甲、乙两地相距680千米，一辆客车和一辆货车同时从两地相对开出，客车每小时行驶75千米，货车每小时行驶65千米，4小时后两车相距多少千米？}

\MethodTitle{方法一}
\Answer{{\Formula{\begin{aligned}[t]&\phantom{={}}680-(75+65)\times4\\&=680-140\times4\\&=680-560\\&=120\text{（千米）}\end{aligned}}}\par 答：两车相距120千米。}

\MethodTitle{方法二：分别求两车的行程后逐次扣除}
\Answer{先算客车和货车各行驶多少千米，再从680千米里扣除。}
\Answer{\Formula{\begin{aligned}[t]&\phantom{={}}680-75\times 4-65\times 4\\
&=680-300-260\\
&=380-260\\
&=120\text{（千米）}\end{aligned}}}
\Answer{答：4小时后两车相距120千米。}
\Problem{③ 玩具厂要生产900个玩具，计划18天完成，实际每天多生产10个，实际多少天完成任务？}

\MethodTitle{方法一}
\Answer{{\Formula{\begin{aligned}[t]&\phantom{={}}900\div(900\div18+10)\\&=900\div(50+10)\\&=900\div60\\&=15\text{（天）}\end{aligned}}}\par 答：实际15天完成任务。}

\MethodTitle{方法二：先求若仍生产18天会多生产多少}
\Answer{计划每天50个，实际每天60个；若仍生产18天，就多生产180个，相当于实际3天的产量，所以可提前3天。}
\Answer{\textbf{计划日产量}\par \Formula{\begin{aligned}[t]&\phantom{={}}900\div 18\\
&=50\text{（个）}\end{aligned}}}
\Answer{\textbf{实际日产量}\par \Formula{\begin{aligned}[t]&\phantom{={}}50+10\\
&=60\text{（个）}\end{aligned}}}
\Answer{\textbf{按18天计算的多产量}\par \Formula{\begin{aligned}[t]&\phantom{={}}10\times 18\\
&=180\text{（个）}\end{aligned}}}
\Answer{\textbf{提前天数}\par \Formula{\begin{aligned}[t]&\phantom{={}}180\div 60\\
&=3\text{（天）}\end{aligned}}}
\Answer{\textbf{实际所用天数}\par \Formula{\begin{aligned}[t]&\phantom{={}}18-3\\
&=15\text{（天）}\end{aligned}}}
\Answer{答：实际15天完成任务。}
\Problem{④ 学校新建的教学楼共有4层，每层有5间教室。学校购进520套双人课桌椅，平均每间教室配多少套？}

\MethodTitle{方法一}
\Answer{{\Formula{\begin{aligned}[t]&\phantom{={}}520\div(5\times4)\\&=520\div20\\&=26\text{（套）}\end{aligned}}}\par 答：平均每间教室配26套。}

\MethodTitle{方法二：先按楼层平均分，再按教室平均分}
\Answer{先把520套分到4层，每层再分到5间教室。}
\Answer{\Formula{\begin{aligned}[t]&\phantom{={}}520\div 4\div 5\\
&=130\div 5\\
&=26\text{（套）}\end{aligned}}}
\Answer{答：平均每间教室配26套。}
\Problem{⑤ 有62名学生，其中会弹钢琴的有11名，会吹竖笛的有56名，两种乐器都不会的有4名。两种乐器都会的有多少名？}

\MethodTitle{方法一}
\Answer{{\Formula{\begin{aligned}[t]&\phantom{={}}11+56-(62-4)\\&=11+56-58\\&=67-58\\&=9\text{（人）}\end{aligned}}}\par 答：两种都会的有9人。}

\MethodTitle{方法二：先求只属于第二类的人数}
\Answer{总人数减去第一类人数，再减去两类都不属于的人数，得到仅属于第二类的人数；再从第二类总人数中扣除。}
\Answer{\Formula{\begin{aligned}[t]&\phantom{={}}56-(62-11-4)\\
&=56-(51-4)\\
&=56-47\\
&=9\text{（人）}\end{aligned}}}
\Answer{答：两种乐器都会的有9人。}
